\documentclass[reprint,prd,amsmath,amssymb,nofootinbib,superscriptaddress]{revtex4-2}

\usepackage{graphicx}
\usepackage{hyperref}
\usepackage{xcolor}
\usepackage{amsmath}
\usepackage{amssymb}
\usepackage{bm}
\usepackage{booktabs}
\usepackage{siunitx}

\hypersetup{colorlinks=true, linkcolor=blue, citecolor=blue, urlcolor=blue}

% --- Custom commands
\newcommand{\jbz}[2]{j_{#1,#2}}      % Bessel zero
\newcommand{\rsqsig}{\sqrt{\sigma}}    % root sigma
\newcommand{\Glueball}[1]{#1}
\newcommand{\Note}[1]{{\color{red}\textbf{[Note: #1]}}}

\begin{document}

\title{On a numerical coincidence between excited \texorpdfstring{$\mathrm{SU}(N)$}{SU(N)} glueball mass ratios at large \texorpdfstring{$N$}{N} and AdS hard-wall Bessel-zero ratios}

\author{K. Remondi\`ere}
\email{kevin.remondiere@gmail.com}
\affiliation{Independent researcher}

\author{Crossed Cosmos collaboration}
\affiliation{ECI/YM research programme}

\date{\today}

\begin{abstract}
We perform a $1/N^2$ cross-extrapolation of the SU($N$) Yang-Mills glueball spectrum data of Athenodorou and Teper (JHEP \textbf{12} (2021) 082) for the assigned $J^{PC}$ channels, comparing the resulting large-$N$ ratios to AdS/QCD hard-wall Bessel-zero predictions. We identify two numerically intriguing coincidences at $N \to \infty$: $m(0^{-+}_{ex_1})/m(0^{++}_{gs}) = 2.288(24)$ matches $\jbz{0}{2}/\jbz{0}{1} = 2.295$ at $0.32\sigma$, and $m(2^{++}_{ex_1})/m(0^{++}_{gs}) = 2.118(13)$ matches $\jbz{2}{1}/\jbz{0}{1} = 2.136$ at $1.34\sigma$. We document the matches' poor cross-$N$ stability (4--12$\sigma$ off at low $N$), the absence of a self-consistent AdS/QCD prediction in the standard hard-wall or soft-wall frameworks, the falsification of the pattern in other excited channels ($0^{++}_{ex_1}$, $2^{-+}_{ex_1}$), and the impact of multi-hypothesis correction (Fisher $p \approx 0.41$ for the joint match, effective $p \gtrsim 1$ after Bonferroni). We list falsifiable predictions for the second-excited 0$^{-+}$ and 2$^{++}$ channels and Sp($2N$) cross-checks. We caveat throughout that the matches do not constitute a discovery but flag them as targets for future high-precision lattice spectroscopy.
\end{abstract}

\maketitle

\section{Introduction}
\label{sec:intro}

Lattice QCD calculations have achieved precise determinations of the SU($N$) Yang-Mills glueball spectrum across a wide range of $N$~\cite{LuciniTeper2010, MorningstarPeardon1999, AthenodorouTeper2021}. The most extensive recent dataset is the Athenodorou-Teper 2021 study~\cite{AthenodorouTeper2021} (henceforth AT 2021), which provides continuum-extrapolated masses of $J^{PC}$ glueballs in units of the string tension $\sqrt{\sigma}$ for $N \in \{2, 3, 4, 5, 6, 8, 10, 12\}$, with the large-$N$ extrapolation derived using canonical 't~Hooft $1/N^2$ corrections.

The AdS/QCD correspondence~\cite{Witten1998, PolchinskiStrassler2000} provides a phenomenological framework predicting hadron spectra from the bulk geometry of a five-dimensional Anti-de~Sitter space. In the hard-wall variant~\cite{ErlichKatzSonStephanov2005, BoschiFilhoBraga2002, DaRoldPomarol2005, deTeramondBrodsky2005}, the spectrum of glueballs and mesons is given by zeros of Bessel functions $J_\nu$. The order $\nu$ depends on the conformal dimension and quantum numbers of the dual boundary operator. Specifically:
\begin{itemize}
  \item For the $\rho$ meson tower (dimension-3 vector operator $\mathrm{Tr}(F_{\mu\nu})$): $m_n^\rho = \jbz{0}{n}/z_{\max}$ with zeros of $J_0$ (Erlich-Katz-Son-Stephanov 2005~\cite{ErlichKatzSonStephanov2005}, eq.~(2.10)).
  \item For the $0^{++}$ scalar glueball (dimension-4 scalar operator $\mathrm{Tr}(F^2)$): $m_n^{0^{++}} = \jbz{2}{n}/z_{\max}$ with zeros of $J_2$ (Boschi-Filho-Braga 2003~\cite{BoschiFilhoBraga2002}).
\end{itemize}
The two towers share the same IR scale $z_{\max}$ in the hard-wall framework (a feature consistent with phenomenological fits)~\cite{ErlichKatzSonStephanov2005, BoschiFilhoBraga2002}.

In a recent cross-$N$ audit of the AT 2021 dataset \cite{HybridFramework2026}, two intriguing numerical coincidences were identified at the SU($N \to \infty$) limit:
\begin{align}
  \frac{m(0^{-+}_{ex_1})}{m(0^{++}_{gs})}\Big|_{N \to \infty} &\approx \frac{\jbz{0}{2}}{\jbz{0}{1}} = 2.295,
  \label{eq:H1prime}\\
  \frac{m(2^{++}_{ex_1})}{m(0^{++}_{gs})}\Big|_{N \to \infty} &\approx \frac{\jbz{2}{1}}{\jbz{0}{1}} = 2.136.
  \label{eq:H2prime}
\end{align}
These two ratios appear to relate the first excited $C=+$ glueballs in the $J^{PC} = 0^{-+}$ and $2^{++}$ channels to specific ratios of Bessel-function zeros that would correspond, in hard-wall AdS/QCD, to first and second zeros of $J_0$, and to the first zero of $J_2$ versus the first zero of $J_0$, respectively.

In this paper we provide an independent verification of these matches, document their poor cross-$N$ stability and theoretical context, identify falsifications of broader claims, and list falsifiable predictions for future lattice work. We caveat throughout that the observed matches do not derive from a self-consistent AdS/QCD framework, and we apply multi-hypothesis statistical correction to report transparently their significance.

\section{Cross-$N$ extrapolation method}
\label{sec:method}

Following the canonical 't~Hooft large-$N$ expansion for $\mathrm{SU}(N)$ Yang-Mills observables, the ratio $r(N) = m_X(N)/m_{0^{++}_{gs}}(N)$ for any glueball channel $X$ admits a $1/N^2$ expansion:
\begin{equation}
  r(N) = c_0 + \frac{c_1}{N^2} + O(1/N^4).
\end{equation}
We fit AT 2021 Tables~34 and~35 data over $N \in \{2,3,4,5,6,8,10,12\}$ using a weighted least-squares procedure, with errors propagated from the published statistical uncertainties on each mass measurement and combined assuming uncorrelated errors. The fit produces best-estimate $c_0$ (asymptotic large-$N$ ratio) and $c_1$ (leading $1/N^2$ correction).

We cross-check by comparing $c_0$ to the published large-$N$ extrapolation in AT 2021 Table~38 (where applicable). In all channels considered, $c_0$ agrees with the Table~38 value within $1.3\sigma$ statistical consistency.

\section{Results}
\label{sec:results}

\subsection{The two key matches}

Independent re-fits with our reproducible code \cite{Code2026} yield:
\begin{align}
  \frac{m(0^{-+}_{ex_1})}{m(0^{++}_{gs})}\Big|_{c_0} &= 2.288(24), \quad c_1 = -0.75(21), \nonumber\\
  &\quad \chi^2/\mathrm{dof} = 1.48, \\
  \frac{m(2^{++}_{ex_1})}{m(0^{++}_{gs})}\Big|_{c_0} &= 2.118(13), \quad c_1 = -0.90(10), \nonumber\\
  &\quad \chi^2/\mathrm{dof} = 0.79.
\end{align}

The discrepancies versus the AdS hard-wall Bessel ratios are:
\begin{align}
  z(0^{-+}_{ex_1}) &= \frac{2.288 - 2.295}{0.024} = -0.32\,\sigma, \\
  z(2^{++}_{ex_1}) &= \frac{2.118 - 2.136}{0.013} = -1.34\,\sigma.
\end{align}

Both are statistically consistent with the AdS hard-wall Bessel targets at less than $1.5\sigma$.

\subsection{Cross-$N$ stability analysis}

Crucially, the matches are not stable across $N$. Table~\ref{tab:crossN} shows the cross-$N$ stability of the same ratios.

\begin{table}[ht]
  \caption{Cross-$N$ stability of the two key ratios. The match with Bessel-zero ratios is poor at low $N$ and emerges only at $N \geq 5$.}
  \label{tab:crossN}
  \begin{tabular}{c|cc|cc}
    \hline
    $N$ & $r(0^{-+}_{ex_1}/0^{++}_{gs})$ & dev$(\jbz{0}{2}/\jbz{0}{1})$ & $r(2^{++}_{ex_1}/0^{++}_{gs})$ & dev$(\jbz{2}{1}/\jbz{0}{1})$ \\
    \hline
    2 & 2.116(42) & $-4.3\sigma$ & 1.910(20) & $-11.5\sigma$ \\
    3 & 2.141(40) & $-3.8\sigma$ & 1.994(17) & $-8.3\sigma$ \\
    4 & 2.241(38) & $-1.4\sigma$ & 2.047(21) & $-4.3\sigma$ \\
    5 & 2.338(42) & $+1.0\sigma$ & 2.094(25) & $-1.7\sigma$ \\
    6 & 2.302(48) & $+0.1\sigma$ & 2.108(31) & $-0.9\sigma$ \\
    8 & 2.226(80) & $-0.9\sigma$ & 2.130(34) & $-0.2\sigma$ \\
    10 & 2.266(53) & $-0.6\sigma$ & 2.118(38) & $-0.5\sigma$ \\
    12 & 2.202(65) & $-1.4\sigma$ & 2.129(29) & $-0.2\sigma$ \\
    \hline
    $\infty$ (Table~38) & 2.295(25) & $-0.0\sigma$ & 2.143(15) & $+0.5\sigma$ \\
    \hline
  \end{tabular}
\end{table}

At $N=2$, the $0^{-+}_{ex_1}/0^{++}_{gs}$ ratio is $-4.3\sigma$ from $\jbz{0}{2}/\jbz{0}{1}$, and the $2^{++}_{ex_1}/0^{++}_{gs}$ ratio is $-11.5\sigma$ off. The pattern of agreement emerges only around $N \geq 5$, suggesting that, if structural, the underlying mechanism is intrinsically a large-$N$ feature.

\subsection{Falsification of broader claims}

We tested whether any general pattern "first excited state of channel $X$ matches Bessel-zero ratio $\jbz{\nu_X}{1}/\jbz{0}{1}$" holds. The results are:

\begin{itemize}
  \item $0^{++}_{ex_1}/0^{++}_{gs}$ at $N \to \infty$ = $1.903(18)$ \emph{vs.} $\jbz{0}{2}/\jbz{0}{1} = 2.295$: $z = -21\sigma$, FALSIFIED.
  \item $2^{-+}_{ex_1}/0^{++}_{gs}$ at $N \to \infty$ = $2.583(21)$ \emph{vs.} $\jbz{2}{1}/\jbz{0}{1} = 2.136$: $z = +21\sigma$, FALSIFIED.
  \item $3^{++}_{gs}/0^{++}_{gs}$ at $N \to \infty$ = $2.364(21)$ \emph{vs.} $\jbz{3}{1}/\jbz{0}{1} = 2.653$: $z = -14\sigma$, FALSIFIED.
\end{itemize}

The two matches in eqs.~(\ref{eq:H1prime},\ref{eq:H2prime}) appear to be \emph{isolated coincidences} in the $0^{-+}_{ex_1}$ and $2^{++}_{ex_1}$ channels, not part of a generalizable pattern.

\subsection{Multi-hypothesis statistical correction}

We tested $\sim 7$ excited-state ratios against $\sim 6$ plausible Bessel-zero ratios, for an effective $\sim 24$ trials. The Fisher-combined $p$-value for the joint $z = -0.32\sigma$ and $z = -1.34\sigma$ observations is
\begin{equation}
  p_{\rm Fisher} = 0.41.
\end{equation}
After Bonferroni correction by a factor $\sim 12$ (half the trial count, accounting for the joint nature), the effective $p$-value is $p_{\rm eff} \gtrsim 1$. The matches are not statistically significant against the null hypothesis of chance coincidence in a multi-trial setting.

\section{Theoretical context}
\label{sec:theory}

\subsection{Standard AdS/QCD hard-wall does not predict the matches}

In the canonical AdS/QCD hard-wall framework~\cite{BoschiFilhoBraga2002, ErlichKatzSonStephanov2005}:
\begin{itemize}
  \item The $0^{++}$ glueball spectrum is $m_n^{0^{++}} = \jbz{2}{n}/z_{\max}$ (zeros of $J_2$).
  \item The $\rho$ meson spectrum is $m_n^{\rho} = \jbz{0}{n}/z_{\max}$ (zeros of $J_0$).
\end{itemize}
The observed ratio $0^{-+}_{ex_1}/0^{++}_{gs} = \jbz{0}{2}/\jbz{0}{1}$ thus relates the lattice glueball quantities to a ratio of consecutive zeros of $J_0$, which (in the standard framework) describes the $\rho$ meson tower internally. There is no canonical hard-wall AdS/QCD reason for the lattice $0^{-+}_{ex_1}/0^{++}_{gs}$ glueball ratio to equal the $\rho$-meson-internal ratio.

Similarly, the observed $2^{++}_{ex_1}/0^{++}_{gs} = \jbz{2}{1}/\jbz{0}{1}$ would (in standard hard-wall) correspond to the lattice ratio of the first excited $2^{++}$ glueball to the ground $0^{++}$ glueball equaling the ratio of the first $J_2$ zero to the first $J_0$ zero. This crosses the meson and glueball towers, again without a derived basis.

Soft-wall variants~\cite{KarchKatzSonStephanov2006, ColangeloDeFazioJugeauNicotri2007} of AdS/QCD predict $m_n^2 \propto n + J/2$ with universal slope, which globally fail the AT 2021 large-$N$ data ($\chi^2/\mathrm{dof} \approx 305$ in joint J^{PC} fit~\cite{HybridFramework2026}).

We conclude that no standard AdS/QCD framework derives the observed matches.

\subsection{Speculative interpretation: phantom ground state}

One could speculate that the lattice $0^{++}_{gs}$ at $m/\sqrt{\sigma} \approx 3.07$ is not the canonical AdS/QCD ground state of any single-trace bulk wavefunction, but rather a "phantom" mode: a multi-glueball threshold, a two-gluon $S$-wave bound state, or a Nambu-Goldstone-like mode of an unrealized broken symmetry. Under this picture, the lattice "first excited state" in $0^{-+}$ and $2^{++}$ channels would be the true holographic ground state of the respective bulk wavefunction tower.

This picture has no derivation in the literature, makes no prediction for which channels are "phantom" (the falsification in $0^{++}_{ex_1}$ and $2^{-+}_{ex_1}$ does not fit), and is speculative. We list it for completeness only.

\section{Falsifiable predictions}
\label{sec:predictions}

If the two matches are structurally meaningful (rather than coincidences), the following further predictions follow:

\begin{enumerate}
  \item Second-excited $0^{-+}$: $r(0^{-+}_{ex_2}/0^{++}_{gs}) \approx \jbz{0}{3}/\jbz{0}{1} = 3.598$, i.e.\ $m(0^{-+}_{ex_2})/\sqrt{\sigma} \approx 11.05$ at $N \to \infty$.
  \item Second-excited $2^{++}$: $r(2^{++}_{ex_2}/0^{++}_{gs}) \approx \jbz{2}{2}/\jbz{0}{1} = 3.500$, i.e.\ $m(2^{++}_{ex_2})/\sqrt{\sigma} \approx 10.75$ at $N \to \infty$.
  \item Sp($2N$) cross-check: under the planar-limit equivalence between $\mathrm{SU}(N)$ and $\mathrm{Sp}(2N)$ gauge theories, the same pattern should hold in the Sp($2N$) lattice spectroscopy. Bennett et al.~\cite{Bennett2020} provide ground and first-excited $0^{++}$ glueball masses in Sp(2$N$) for $N=1,\ldots,4$ with large-$N$ extrapolation. The corresponding Sp($2N$) ratio $m(0^{++}_{ex_1})/m(0^{++}_{gs})$ can be directly compared to the SU($N$) value 1.903(18). To test the $0^{-+}_{ex_1}$ and $2^{++}_{ex_1}$ predictions specifically, additional Sp($2N$) excited-state spectroscopy in those channels would be required. A confirmation at the SU/Sp planar limit would falsify the "$\mathrm{SU}(N)$ accidental coincidence" alternative.
\end{enumerate}

Current AT 2021 data are too sparse on $ex_2$ excited states to test predictions (1) and (2) at $N \to \infty$.

\section{Conclusions}
\label{sec:conclusions}

We have independently verified that two specific numerical coincidences hold at the SU($N \to \infty$) limit of the AT 2021 glueball spectrum data, with statistical significance of $0.32\sigma$ and $1.34\sigma$ in the $0^{-+}_{ex_1}/0^{++}_{gs}$ and $2^{++}_{ex_1}/0^{++}_{gs}$ channels respectively, against AdS hard-wall Bessel-zero ratios $\jbz{0}{2}/\jbz{0}{1}$ and $\jbz{2}{1}/\jbz{0}{1}$.

We have also documented:
\begin{itemize}
  \item Poor cross-$N$ stability: the matches are $4-12\sigma$ off at low $N$ and emerge only at $N \geq 5$.
  \item Absence of a self-consistent AdS/QCD framework predicting the matches.
  \item Falsification of any general "EX1 ↔ ground-Bessel" pattern in other excited channels.
  \item Multi-hypothesis correction reducing Fisher $p \approx 0.41$ to effective $p \gtrsim 1$.
\end{itemize}

The matches do not constitute a discovery but they are intriguing and warrant inclusion in any future cross-$N$ analysis of the SU($N$) glueball spectrum. We list specific falsifiable predictions for the $0^{-+}_{ex_2}$, $2^{++}_{ex_2}$, and Sp($2N$) lattice data, which could either elevate the present observation to a structural feature or definitively reject it as a coincidence.

\section*{Acknowledgments}

The cross-$N$ extrapolation reproducible code is available at \cite{Code2026}.

\begin{thebibliography}{99}

\bibitem{Witten1998}
  E.~Witten,
  ``Anti-de Sitter Space, Thermal Phase Transition, and Confinement in Gauge Theories,''
  Adv.~Theor.~Math.~Phys.\ \textbf{2} (1998) 505;
  \href{https://arxiv.org/abs/hep-th/9803131}{arXiv:hep-th/9803131}.

\bibitem{PolchinskiStrassler2000}
  J.~Polchinski and M.~J.~Strassler,
  ``The String Dual of a Confining Four-Dimensional Gauge Theory,''
  \href{https://arxiv.org/abs/hep-th/0003136}{arXiv:hep-th/0003136} (2000).

\bibitem{ErlichKatzSonStephanov2005}
  J.~Erlich, E.~Katz, D.~T.~Son and M.~A.~Stephanov,
  ``QCD and a Holographic Model of Hadrons,''
  Phys.~Rev.~Lett.\ \textbf{95}, 261602 (2005);
  \href{https://arxiv.org/abs/hep-ph/0501128}{arXiv:hep-ph/0501128}.

\bibitem{BoschiFilhoBraga2002}
  H.~Boschi-Filho and N.~R.~F.~Braga,
  ``QCD/String holographic mapping and glueball mass spectrum,''
  Eur.\ Phys.\ J.\ C \textbf{32} (2003) 529;
  \href{https://arxiv.org/abs/hep-th/0209080}{arXiv:hep-th/0209080}.

\bibitem{DaRoldPomarol2005}
  L.~Da Rold and A.~Pomarol,
  ``Chiral symmetry breaking from five dimensional spaces,''
  Nucl.\ Phys.\ B \textbf{721}, 79 (2005);
  \href{https://arxiv.org/abs/hep-ph/0501218}{arXiv:hep-ph/0501218}.

\bibitem{deTeramondBrodsky2005}
  G.~F.~de Teramond and S.~J.~Brodsky,
  ``Hadronic Spectrum of a Holographic Dual of QCD,''
  Phys.\ Rev.\ Lett.\ \textbf{94}, 201601 (2005);
  \href{https://arxiv.org/abs/hep-th/0501022}{arXiv:hep-th/0501022}.

\bibitem{KarchKatzSonStephanov2006}
  A.~Karch, E.~Katz, D.~T.~Son and M.~A.~Stephanov,
  ``Linear Confinement and AdS/QCD,''
  Phys.\ Rev.\ D \textbf{74}, 015005 (2006);
  \href{https://arxiv.org/abs/hep-ph/0602229}{arXiv:hep-ph/0602229}.

\bibitem{ColangeloDeFazioJugeauNicotri2007}
  P.~Colangelo, F.~De Fazio, F.~Jugeau and S.~Nicotri,
  ``On the light glueball spectrum in a holographic description of QCD,''
  Phys.\ Lett.\ B \textbf{652}, 73 (2007);
  \href{https://arxiv.org/abs/hep-ph/0703316}{arXiv:hep-ph/0703316}.

\bibitem{MorningstarPeardon1999}
  C.~J.~Morningstar and M.~J.~Peardon,
  ``The glueball spectrum from an anisotropic lattice study,''
  Phys.\ Rev.\ D \textbf{60}, 034509 (1999).

\bibitem{LuciniTeper2010}
  B.~Lucini, A.~Rago and E.~Rinaldi,
  ``Glueball masses in the large $N$ limit,''
  JHEP \textbf{08}, 119 (2010).

\bibitem{AthenodorouTeper2021}
  A.~Athenodorou and M.~Teper,
  ``SU($N$) gauge theories in 3+1 dimensions: glueball spectrum, string tensions and topology,''
  JHEP \textbf{12} (2021) 082;
  \href{https://arxiv.org/abs/2106.00364}{arXiv:2106.00364}.

\bibitem{Bennett2020}
  E.~Bennett, J.~Holligan, D.~K.~Hong, J.~W.~Lee, C.-J.~D.~Lin, B.~Lucini, M.~Piai and D.~Vadacchino,
  ``Glueballs and Strings in Sp($2N$) Yang-Mills theories,''
  Phys.\ Rev.\ D \textbf{103}, 054509 (2021);
  \href{https://arxiv.org/abs/2010.15781}{arXiv:2010.15781}.

\bibitem{Shanahan2025}
  R.~Abbott, D.~C.~Hackett, D.~A.~Pefkou, F.~Romero-L\'opez, and P.~E.~Shanahan,
  ``Lattice evidence that scalar glueballs are small,''
  Phys.\ Rev.\ Lett.\ \textbf{136}, 041901 (2026);
  \href{https://arxiv.org/abs/2508.21821}{arXiv:2508.21821}.

\bibitem{HybridFramework2026}
  Opus 4.7 max-effort hybrid framework deep-dive (internal report),
  \texttt{HYBRID\_FRAMEWORK\_2026-05-15.md}, May 15, 2026.

\bibitem{Code2026}
  Reproducible Python scripts for cross-$N$ fits and falsifiable predictions:
  \texttt{paper\_6prime\_cross\_N\_fits.py} and \texttt{paper\_6prime\_bessel\_predictions.py},
  May 2026.

\end{thebibliography}

\end{document}
